% Complete isolated algebraic candidate. No geometric gauge is asserted here.
\subsection{The Fermat cubic does not divide a nonzero Moore determinant}

Keep the characteristic-two Moore matrix (DE-MC1), with
\(G=\mathbf F_3^2\), lexicographic coordinate order, and algebraically
closed ground field \(k\). Put
\[
 \mathcal F(x)=\sum_{g\in G}x_g^3,\qquad
 R_c(x)=\operatorname{rem}_{x_0}
       \bigl(D_c(x),x_0^3+\sum_{g\ne0}x_g^3\bigr).
 \tag{DE-MF1}
\]
Division is monic, so it commutes with specialization in \(c\), and
\(\mathcal F\mid D_c\) is equivalent to \(R_c=0\).

\begin{lemma}[Exact coefficients of the Fermat remainder]
\label{lem:DE-Fermat-coefficient-reductions}
\status{Candidate; complete proof and standalone exact certificate supplied.}
If \(R_c=0\), then
\[
 \sum_g c_g^3=0,\qquad
 P_H(c):=\sum_{r\in G/H}\prod_{g\in r+H}c_g=0
       \quad\text{for each line }H\subset G.
 \tag{DE-MF2}
\]
For \(c\ne0\), the only possible supports are the four types in
(DE-MC3), or the full support \(G\). On the two-, four-, and
six-coordinate representatives, all eight coefficients
(DE-MC4)--(DE-MC6) are also actual coefficients of \(R_c\).
On the eight-coordinate representative, once
\(c_v^3=c_{-v}^3\) is known, the only phase patterns (DE-MC7)
consistent with \(R_c=0\) are the nine characters of \(G\).
\end{lemma}
\begin{proof}
The exact universal reconstruction uses the exponent counts in the
proof of Lemma~\ref{lem:DE-Moore-support}. After cancellation of the
\(9!\) permutation terms, it has \(55{,}296\) monomials and
\(2{,}080\) coefficient polynomials. For every monomial write
its \(x_0\)-exponent as \(3q+r\), \(0\le r<3\), and replace
\(x_0^{3q+r}\) by \(x_0^r(\sum_{g\ne0}x_g^3)^q\).
Cancel equal monomials in characteristic two. The resulting
\(R_c\) has \(1{,}085\) coefficient polynomials and
\(47{,}349\) monomials.

The coefficient of \(\prod_gx_g\) is
\((\sum_gc_g^3)^3\). It is unchanged by monic division, since
every nonconstant contribution from the replacement has an
\(x_g\)-exponent at least three. For each of the four line
directions, the cube \(P_H(c)^3\) is an actual remainder
coefficient. The standalone certificate reconstructs these five
identities over \(\mathbf F_2[c]\), retaining the exponent
locator for each. A field has no nonzero nilpotents, so these
identities prove (DE-MF2).

For each of the \(511\) nonempty supports, discard terms using
a coordinate outside that support. Exactly \(399\) supports
have a remainder coefficient consisting of one nonzero monomial;
its exponent vectors are retained. The other \(112\) are exactly
the \(36,54,12,9\) supports of sizes \(2,4,6,8\) in (DE-MC3),
and the full support. All \(432\) affine index maps are checked,
and a map from the indicated representative is supplied for
every surviving support. The eight factored polynomials
(DE-MC4)--(DE-MC6) are reconstructed and located in these restricted
remainders, with no division by a possibly zero coordinate.

Finally substitute (DE-MC7) into the entire eight-support
remainder, retaining the four independent nonzero variables
\(d_v\). Of the \(81\) root-of-unity choices, exactly the
nine character choices give zero. For each other choice the
certificate retains a coefficient equal to a nonzero
\(\mathbf F_4\)-scalar times a monomial in the \(d_v\).
This proves the last assertion over \(k\), rather than by
sampling parameter values. The opposite-cube equalities are
not inferred from the old coefficient (DE-MC9): that coefficient
changes under division. They are proved below by a separate
line-configuration argument.
\end{proof}

\begin{lemma}[Affine-line restrictions and their two configurations]
\label{lem:DE-Fermat-affine-lines}
\status{Candidate; complete proof and exact orbit checks supplied.}
If \(\mathcal F\mid D_c\), then in each parallel class of
affine lines in \(G\), at least one line \(J\) satisfies
\[
                         \sum_{g\in J}c_g^3=0.       \tag{DE-MF3}
\]
Choosing one such line in each of the four directions gives
\(81\) configurations: nine concurrent configurations and
seventy-two configurations with a unique triple point and
the fourth line off that point. The affine group is transitive
on each class.
\end{lemma}
\begin{proof}
Restrict the independent variables \(x\) to a line subgroup
\(H=\langle v\rangle\). Group rows as \(r+av\) and columns
as \(-r+bv\), where \(a,b\in\mathbf F_3\) and \(r\) runs
through representatives of \(G/H\). Off these three blocks
the matrix is zero, and a block has entries
\(c_{2r+(a-b)v}x_{(a+b)v}\). The three-by-three determinant is
\[
 \Bigl(\prod_{g\in2r+H}c_g\Bigr)\sum_{h\in H}x_h^3
 +\Bigl(\sum_{g\in2r+H}c_g^3\Bigr)\prod_{h\in H}x_h.
 \tag{DE-MF4}
\]
The full restricted determinant is the product of these three
expressions. Modulo \(\sum_{h\in H}x_h^3\), it becomes the
product of the three cube sums times \((\prod_{h\in H}x_h)^3\).
The Fermat polynomial does not divide that coordinate monomial.
Thus one cube sum vanishes, proving (DE-MF3).

Translate the intersection of two chosen lines to zero and
make them the coordinate axes. The other two have equations
\(x+y=\alpha\), \(x-y=\beta\). If both constants vanish,
all four lines concur. Otherwise a triple point exists:
this is immediate if one constant vanishes, and if neither
does, \(\beta=\alpha\) or \(\beta=-\alpha\) in \(\mathbf F_3\).
Two different triple points would force all four lines through
their common intersection, so the triple point is unique.
There are nine concurrent choices, and
\(9\cdot4\cdot2=72\) choices of a triple point, omitted
direction and fourth line. These are the two affine orbits.

The covariances used here preserve Fermat divisibility. For
\(c'_g=c_{Ag+b}\), \(x'_g=x_{Ag}\), the matrix
\(M_{c'}(x')\) is \(M_c(x)\) with row index \(Aa+2b\)
and column index \(Ab+b\); sums and differences give the
claimed entries. Both index maps are permutations and
\(\mathcal F(x')=\mathcal F(x)\). Multiplying \(c_g\) by
\(\zeta^{\ell(g)}\) is a row and column scaling of the Moore
matrix and preserves each cube sum and each \(P_H\).
Homogeneous rescaling of \(c\) also preserves every zero
condition. These identities justify all normalizations below.
\end{proof}

\begin{lemma}[Full support is impossible; the eight-support cubes agree]
\label{lem:DE-Fermat-full-and-eight}
\status{Candidate; complete zero-orbit and polynomial-identity proof supplied.}
Suppose \(c\ne0\) and \(\mathcal F\mid D_c\).
The support cannot be all of \(G\). If it is
\(G\setminus\{h\}\), then
\(c_{h+v}^3=c_{h-v}^3\) for every \(v\).
\end{lemma}
\begin{proof}
Choose the four lines in Lemma~\ref{lem:DE-Fermat-affine-lines}.
If they concur at \(h\), summing their four cube-sum equations
gives \(\sum_gc_g^3=c_h^3\). By (DE-MF2), \(c_h=0\).
This excludes full support; on eight-support the center is
the unique missing point, and each of its four lines gives
the required opposite-cube equality.

In the other orbit take the lines
\(x=0,y=0,x+y=0,x-y=1\). In lexicographic order write
\[
 c=(e,a,u,b,d,v,w,t,f),\qquad E=e^3, A=a^3, B=b^3.
\]
Their four equations and \(\sum c_g^3=0\) give
\[
 u^3=E+A,\quad w^3=E+B,\quad v^3=A+B,\quad
 t^3=E+A+B,\quad f^3=d^3.                           \tag{DE-MF5}
\]
When \(d,f\ne0\), a character scaling makes \(f=d\):
their quotient is in \(\mu_3\), and a linear character can
take any prescribed value on \((2,2)-(1,1)\).
Three of the four product equations (DE-MF2) then become
\[
 eau+d(bv+wt)=0,\quad ebw+d(at+uv)=0,\quad
 evt+d(ab+uw)=0.                                    \tag{DE-MF6}
\]

For full support, scale \(e=1\), put
\(U=1+A,W=1+B,V=A+B,T=1+A+B,D=d^3\), and set
\(K=d^2abuvwt\ne0\). The cube relations give
\[
 BV+WT=U,\qquad AT+UV=B,\qquad AB+UW=T.
\]
Cubing the three equalities in (DE-MF6), and using
\((r+s)^3=r^3+s^3+rs(r+s)\), gives respectively
\[
             (D+A)U=K,\qquad(D+W)B=K,\qquad(D+V)T=K.
 \tag{DE-MF7}
\]
But \((D+A)U+(D+W)B=(D+V)T\). The first two equations
make its left side zero, and the third makes its right side
\(K\ne0\), a contradiction. The certificate also verifies
an explicit ideal identity for this contradiction, including
all cross terms in cubing the three equations.

Now assume exactly one coordinate is zero. The stabilizer of
the selected four lines fixes their triple point. In a basis
\(v=(1,1),q=(1,0)\), it consists of
\(v\mapsto\epsilon v,q\mapsto q+\tau v\), where
\(\epsilon\in\{1,-1\}\), \(\tau\in\mathbf F_3\).
Its four point orbits are
\[
 \{0\},\qquad\{2,3,7\},\qquad\{1,5,6\},\qquad\{4,8\},
 \tag{DE-MF8}
\]
using indices \(0,\ldots,8\). They are, respectively, the
triple point, the three double points, the three single
points and the two uncovered points. The six transformations
and all four orbits are checked exactly in the certificate.

If the zero is uncovered, (DE-MF5) forces both \(d=f=0\),
contrary to uniqueness. If the zero is the triple point
\(e=0\), (DE-MF5) immediately gives all four opposite-cube
equalities about zero. In the remaining two orbits \(e,d,f\)
are nonzero, so normalize \(e=1,f=d\).

For a single-point zero take \(a=0\). Then
\(U=1,V=B,W=T=1+B\). The first equation of (DE-MF6)
gives \(bv+wt=0\). Its cube is
\(BV+WT+bvwt(bv+wt)=B^2+(1+B)^2=1\), a contradiction.
For a double-point zero take \(u=0\). Then
\(A=1,T=B,W=V=1+B\). The second equation of (DE-MF6)
is \(bw=dat\); cubing gives \(BW=DB\). Since \(b\ne0\),
\(D=W=V\). Thus the cube pairs about the missing point
\((0,2)\), with index2, agree: the pairs are
\((0,1),(3,7),(4,6),(5,8)\), with common cubes
\(1,B,1+B,1+B\), respectively. Covariance transports this
conclusion to every double-point zero. All cases are exhausted.
\end{proof}

\begin{theorem}[Exact Fermat nondivisibility]
\label{thm:DE-Moore-Fermat-nondivisibility}
\status{Candidate; complete proof and exact certificate, independent audits pending.}
For every \(c\in k^9\),
\[
       \mathcal F(x)\mid D_c(x)
             \quad\Longleftrightarrow\quad D_c(x)\equiv0.
 \tag{DE-MF9}
\]
For \(c\ne0\), this is equivalent to membership in one of
the eighty-one spaces \(V_{h,\ell}\) of (DE-MC2).
\end{theorem}
\begin{proof}
The reverse implication is immediate, as is the case \(c=0\).
Assume \(c\ne0\) and \(\mathcal F\mid D_c\).
Lemma~\ref{lem:DE-Fermat-coefficient-reductions} reduces its
support to the five listed types, and
Lemma~\ref{lem:DE-Fermat-full-and-eight} excludes full support.
On two- and four-coordinate supports, the actual remainder
coefficients (DE-MC4)--(DE-MC5) give the opposite cube ratios,
which extend to a character exactly as in the proof of
Theorem~\ref{thm:DE-Moore-singular-parameters}.
On six-coordinate support, the five actual coefficients
(DE-MC6) give the sum, product and three cyclic matching
conditions. Apply the proved
Lemma~\ref{lem:DE-Moore-six-matching}, and then the product
identity makes the three cube-root phases affine linear.
The full reflection-and-character construction in that theorem
therefore again gives \(c\in V_{h,\ell}\).

On eight-coordinate support the preceding lemma supplies the
previously missing opposite-cube equalities. Translate the
missing coordinate to zero and apply the independent-variable
phase assertion of
Lemma~\ref{lem:DE-Fermat-coefficient-reductions}. Only a
character phase remains, giving the same membership.
Finally every such Moore matrix is made alternating of odd
size by the explicit row scaling in
Theorem~\ref{thm:DE-Moore-singular-parameters}; its determinant
is identically zero. This proves both equivalences.
\end{proof}

\begin{remark}[Certificate, correction history and precise scope]
\label{rem:DE-Moore-Fermat-scope}
\status{Candidate scope; no untyped geometric or monodromy implication.}
The standalone program
\path{code/obstructions_crystals/verify_Moore_Fermat_nondivisibility.py}
uses no external input or imported output. Its completed author
run passed at 02:34:06.439903 UTC on 20 November 2026 in
\(4.910654410996358\) seconds. It reconstructs the determinant
and monic remainder, retains all \(399\) support exclusions,
checks all \(432\) affine maps, the eight old factors and four
parallel-line product cubes, all \(81\) line choices and the
six-element stabilizer, the zero-position algebra, the explicit
full-torus ideal identity, and all \(81\) eight-support phases.
Every noncharacter phase has a retained nonzero monomial witness.

The earlier exploratory search found no opposite-pair binomial
among the eight-support remainder coefficients. Its four empty
lists are retained; no such binomial is used here. The proof
instead establishes the cube ratios by (DE-MF5)--(DE-MF8).
An initial execution of the new checker rejected a proposed
ideal expression that omitted the \(rs(r+s)\) term in a cubed
sum. Its exact source and standard streams are retained; the
corrected ideal expression includes the corresponding multiples
of the three equations and passes. The cubed-equality argument
(DE-MF7) itself does not omit those terms. Earlier finite-field
probes and the counterexample to sufficiency of the affine-line
conditions alone are also preserved; they are not used to infer
the universal theorem.

This is a set-theoretic universal matrix statement, not a
scheme-theoretic radical computation. To apply it to a Cartier
wall one must still identify every relevant determinant and the
degree-one wall with this same \(x\)-coordinate Fermat cubic;
separate arbitrary linear coordinate changes would not suffice.
No new global monodromy generation or geometric conclusion is
asserted by this subsection alone.
\end{remark}
